\section{One Sacrifice, One Priest, One Worshipping Body}

Before the Roman books are compared, the thing enacted must be confessed. At the Last Supper Christ takes bread and wine, gives thanks, identifies them with his Body given and Blood poured out, and commands the apostolic memorial. On Calvary he offers himself once for all. In the Resurrection and Ascension his priesthood does not become past: the Lamb who was slain lives and intercedes. The Mass is therefore neither another Calvary nor a congregational recollection of an absent event. It is the sacramental making-present of the one sacrifice under signs instituted by Christ (Luke 22:14--20; 1~Cor. 10:16--21; 11:23--29; Heb. 7--10; Rev. 5).

Trent's Session XXII teaches that the same Christ is victim and principal priest in Cross and Mass; only the manner of offering differs. The unbloody sacramental offering represents, makes present, and applies the bloody sacrifice without repeating it. It is truly propitiatory, offered for living and dead, and entrusted to the ministerial priesthood. Vatican II speaks the same faith: at the Last Supper Christ instituted the Eucharistic sacrifice “in order to perpetuate the sacrifice of the Cross throughout the centuries” and entrusted to the Church the memorial of his death and resurrection (\emph{Sacrosanctum Concilium} 47). The postconciliar General Instruction places its own account explicitly in continuity with Trent.

\subsection{Christ's action and the Church's action}

The celebrant acts \latin{in persona Christi Capitis}; the baptized exercise their common priesthood by joining themselves to Christ's offering and offering themselves. These modes are ordered but not interchangeable. \emph{Lumen gentium} 10 says the ministerial and common priesthood differ in essence and not only degree. A priest does not acquire a larger share of a generic congregational power; sacramental Orders configures him to act as Christ the Head in a mode the unordained cannot supply. The people are not spectators for that reason. Their Amen, prayer, chant, silence, bodily discipline, sacrifice of praise, and self-offering are real liturgical acts precisely because Christ gathers a hierarchically ordered Body.

Aquinas supplies the causal grammar. Christ is the principal minister; the ordained priest is his living instrument; the sacramental signs are instruments by which the Passion reaches persons in time (ST III, qq.~62, 73--83). The minister's sanctity affects the moral quality and edification of his act but not Christ's fidelity to a valid sacrament. This is the anti-Donatist truth Augustine condensed by saying that whether Peter or Judas baptizes, Christ baptizes. It also explains why sacramental objectivity never licenses casual celebration: an instrument acts truly, and sacrilege is more grave because Christ remains faithful.

\subsection{Rite, substance, and ecclesial authority}

Christ determines sacramental substance; the Church orders its liturgical clothing and safeguards the sign. Trent Session VII, canon 13 rejects contempt for the Church's received and approved rites or private alteration of them. Pius XII's \emph{Sacramentum ordinis} distinguishes the substance Christ instituted from elements the Church can determine. Canon 841 reserves to supreme ecclesiastical authority what validity requires; canon 846 requires ministers to follow approved books faithfully and forbids private addition, omission, or alteration.

This authority is real but ministerial. The Church cannot decree that beer is Eucharistic wine, an unordained person a sacramental priest, or a non-Trinitarian washing Christian Baptism. She can determine an approved consecratory prayer, calendar, ceremonial, lectionary, language, posture, number of anointings, and disciplinary conditions. A pope's authority to promulgate the 1570 Missal and his successors' authority to revise it arise from the same office. “Pius V could bind every future pope never to revise his disciplinary book” is not a defense of papacy but a contradiction of its continuing legislative competence.

\subsection{The indefectibility boundary}

The Church cannot give her faithful, as the normative Eucharistic rite of the Roman Church, a sacrament that is intrinsically false worship or incapable of validity. This does not mean every reform is maximally prudent, every translation beautiful, or every disciplinary permission equally conducive to reverence. It means that theories declaring the officially promulgated postconciliar Order intrinsically invalid or heretical collide with the Church's indefectibility and her authoritative sacramental judgment.

The reverse mistake treats authorization as historical perfection. The Roman Mass did not begin when authority printed it, and a lawful reconstruction can still be evaluated against the authority's stated principles. \emph{Sacrosanctum Concilium} 23 makes careful theological, historical, and pastoral investigation, genuine and certain need, and organic growth from existing forms the reform's own criteria. Applying those criteria is not private jurisdiction over the liturgy. It is disciplined reception of the council's instruction.

\subsection{What ritual form can emphasize}

Every orthodox Mass contains more reality than any finite set of words can exhaust. A rite distributes emphases by repetition, silence, gesture, sequence, and vocabulary. Propitiation can be doctrinally present throughout one Missal yet receive more repeated verbal emphasis in another. The ecclesial assembly can be real in every Low Mass yet appear more visibly in dialogical rites. Scripture can permeate the older Ordinary while a much greater range is proclaimed in the newer lectionary. The Roman Canon can remain available in the postconciliar Missal while weekly options make its practical dominance disappear.

Ritual therefore forms perception without mechanically determining belief. Lex orandi and lex credendi do not mean that every omitted phrase deletes a dogma or that a single prayer proves the average communicant's catechesis. They mean that the Church's worship expresses, guards, and habituates faith. Comparative theology asks what a complete ritual ecology repeatedly makes easy to perceive, not whether either book can be forced into a heretical slogan.

\begin{studybox}{The sacramental center}
At either authorized altar, a validly ordained priest intends the Church's act over valid bread and wine and pronounces the approved Eucharistic Prayer. Christ acts. The one sacrifice becomes sacramentally present; bread and wine are transubstantiated; the Church offers through and with her Head; the faithful are ordered to sacrificial Communion and charity. No language, orientation, musical style, or personal preference creates that reality. All of them can help or hinder its fitting human reception.
\end{studybox}
